\documentclass[12pt]{article}
\usepackage[utf8]{inputenc}
\usepackage[T1]{fontenc}
\usepackage{hyperref}
\title{Yawanawa Language Resources for K-12 Curriculum Development}
\author{Compiled by Claude Sonnet 5.5 (claude-sonnet-5-5) for Andr\'as Kornai}
\date{}
\begin{document}
\maketitle

\section{Introduction}
Yawanawa (ISO 639-3 ywn; Glottocode yawa1260) is a Panoan language of Acre, Brazil, spoken by about 1,000 people. The searches found some material in the language and about it. Literacy is reported as very low, and the school resources found are mostly teaching guides and textbooks. The searches found no dedicated Yawanawa dictionary.

\section{Primary Educational Resources}
\begin{itemize}
\item \emph{Yawanawahãu Tãnãti: Nuk\~e Matematica}, a Yawanawa-language mathematics book. A numerals database cites it without a publisher or date. The searches did not find an online copy. Citation page: \url{https://lingweb.eva.mpg.de/channumerals/Yawanawa.htm}
\item \emph{Caderno Pedagógico de Biologia Yawanawá}, a first-year high-school biology guide for Yawanawá schools in the Tarauacá area. It was produced in a professional master's programme at the Universidade Federal do Acre. The searches found it only in a summary, and this report gives no URL for it.
\end{itemize}

\section{Linguistic and Reference Materials}
\begin{itemize}
\item Grambank entry for Yawanawa. It cites Aldir Santos de Paula, \emph{A língua dos índios Yawanawá do Acre} (Maceió, 2007, 163 pp.), and a 2013 thesis by de Souza on nominal expressions. \url{https://grambank.clld.org/languages/yawa1260}
\item Channumerals database, Yawanawa numeral data: \url{https://lingweb.eva.mpg.de/channumerals/Yawanawa.htm}
\end{itemize}

\section{Religious and Community Resources}
\begin{itemize}
\item \emph{Costumes e Tradições do Povo Yawanawá} (2006), produced by the Organização dos Professores Indígenas do Acre with the Comissão Pró-Índio do Acre. It is a PDF. \url{https://cpiacre.org.br/wp-content/uploads/2020/03/Costumes-e-Tradicoes-do-Povo-Yawanawa.pdf}
\end{itemize}
No scripture translations were found.

\section{Digital and Audio Resources}
The summary of the biology guide mentions existing Yawanawá films, CDs and websites. No specific audio or video was identified with a URL.

\section{Curriculum Development Context}
The materials above suggest some bilingual school activity in Yawanawá schools in Acre. Whether it is systematic, or how far it uses the language, could not be established from the search results. The OPIAC and CPI-Acre publication and the university-based pedagogical work are the natural starting points. Further searches, for example for a ``dicionário Yawanawá'' and the Universidade Federal do Acre repositories, would be needed.

\end{document}